\begin{figure}[ht]
\centering
\begin{tikzpicture}[x=1cm, y=1cm, font=\small]

% LEFT PANEL: active constraint. Minimise (1/2)||x||^2 subject to
% a^T x >= b with b > 0 so the origin is infeasible; optimum on boundary.
\begin{scope}[shift={(0,0)}]
  \node[anchor=south, font=\small] at (0, 3.2) {constraint active};
  \draw[->, thick] (-3.0, 0) -- (3.0, 0) node[right, font=\scriptsize] {$x_{1}$};
  \draw[->, thick] (0, -2.4) -- (0, 2.9) node[above, font=\scriptsize] {$x_{2}$};
  % Objective level circles
  \foreach \r in {0.5, 1.0, 1.5} { \draw[engblue, thin] (0,0) circle (\r); }
  % Constraint line a^T x = b, a = (1,1), b = 1.4: x_2 = 1.4 - x_1
  % Feasible side a^T x >= b is the upper-right half-plane.
  \draw[engred, very thick] (-1.3, 2.7) -- (2.7, -1.3);
  % Shade feasible half-plane lightly
  \fill[engred!8] (-1.3, 2.7) -- (2.7, -1.3) -- (2.9, 2.7) -- cycle;
  \node[engred, anchor=west, font=\scriptsize] at (1.6, 0.5)
    {$\vect{a}^{\transpose}\vect{x} = b$};
  % Optimum = projection of origin onto line = (b/||a||^2) a = (0.7,0.7)
  \fill[engblack] (0.7, 0.7) circle (2.2pt);
  \node[engblack, anchor=south west, font=\scriptsize] at (0.78, 0.74)
    {$\vect{x}^{*}$};
  % gradient of f points outward along x^*, lambda* grad f1 = -lambda* a inward
  \draw[->, engblue, thick] (0.7, 0.7) -- (1.2, 1.2);
  \node[engblue, anchor=west, font=\scriptsize] at (1.2, 1.25)
    {$\grad f$};
  \fill[engblack] (0,0) circle (1.3pt);
  \node[anchor=north, font=\scriptsize] at (0, -2.05) {$\lambda^{*} > 0$};
\end{scope}

% RIGHT PANEL: inactive constraint. Same objective but b < 0 so the
% unconstrained optimum (origin) is already feasible.
\begin{scope}[shift={(7.5,0)}]
  \node[anchor=south, font=\small] at (0, 3.2) {constraint inactive};
  \draw[->, thick] (-3.0, 0) -- (3.0, 0) node[right, font=\scriptsize] {$x_{1}$};
  \draw[->, thick] (0, -2.4) -- (0, 2.9) node[above, font=\scriptsize] {$x_{2}$};
  \foreach \r in {0.5, 1.0, 1.5} { \draw[engblue, thin] (0,0) circle (\r); }
  % Constraint line pushed away from the origin: a^T x = -1.2,
  % feasible side a^T x >= -1.2 contains the origin.
  \draw[engred, very thick] (-2.8, 1.6) -- (1.6, -2.8);
  \fill[engred!8] (-2.8, 1.6) -- (1.6, -2.8) -- (2.9, 2.9) -- (-2.8, 2.9) -- cycle;
  \node[engred, anchor=west, font=\scriptsize] at (-2.7, 0.9)
    {$\vect{a}^{\transpose}\vect{x} = b$};
  % Optimum = unconstrained minimum at the origin
  \fill[engblack] (0,0) circle (2.2pt);
  \node[engblack, anchor=south west, font=\scriptsize] at (0.08, 0.08)
    {$\vect{x}^{*}=\vect{0}$};
  \node[anchor=north, font=\scriptsize] at (0, -2.05) {$\lambda^{*} = 0$};
\end{scope}

\end{tikzpicture}
\caption[KKT geometry: active and inactive inequality
constraints]{The two cases of complementary slackness for the
projection problem $\min \tfrac{1}{2}\norm{\vect{x}}^{2}$ subject
to $\vect{a}^{\transpose}\vect{x} \geq b$. When $b > 0$ (left) the
unconstrained optimum (the origin) is infeasible, the constraint
binds, the optimum sits on the boundary hyperplane, and the
multiplier $\lambda^{*} = b/\norm{\vect{a}}^{2}$ is strictly
positive; the objective gradient at $\vect{x}^{*}$ is parallel to
the constraint normal. When $b < 0$ (right) the origin is already
feasible, the constraint is slack, the optimum is the
unconstrained minimum, and the multiplier $\lambda^{*} = 0$.
Complementary slackness $\lambda^{*}(b - \vect{a}^{\transpose}
\vect{x}^{*}) = 0$ holds in both cases: one factor is always zero.}
\label{fig:vol02:ch15:kkt-geometry}
\end{figure}
